%! The Language of Variation — Unit 1, Lesson 1.1 — Warm-Up (Answer Key)
% Mirrors the blank exactly apart from the package swap, the pageheader, and the filled
% answers. ONE PAGE, like the blank. NO teachernote here — teacher prose lives in the plan.
\documentclass[10pt]{article}
\usepackage{apstats-article}
\usepackage{apstats-key}   % pulls in -boxes; do NOT also load -boxes

\begin{document}

\pageheader{Unit 1, Lesson 1.1}{Warm-Up --- Answer Key}

\vspace{0.15in}

\boxguard[24]
\begin{hookbox}
\small
A music teacher keeps this sheet on every student who rents an instrument.

\vspace{6pt}
\begin{center}
\renewcommand{\arraystretch}{1.25}
\begin{tabular}{@{} l l c c @{}}
\rowcolor{sky}
\textbf{Student} & \textbf{Instrument} & \textbf{Weeks rented} & \textbf{Locker \#} \\
\hline
Ibarra & Clarinet & 14 & 208 \\
Okafor & Trumpet  &  6 & 145 \\
Petrov & Clarinet & 22 & 331 \\
\hline
\end{tabular}
\end{center}
\vspace{2pt}
\end{hookbox}

\vspace{0.14in}

\begin{enumerate}[leftmargin=1.7em, itemsep=0.15in]

  \item In this table, the \textbf{individuals} are \ans{the students renting instruments}.

        Name one \textbf{variable} \ans{Instrument} and one of its \textbf{values}
        \ans{Clarinet}.

  \item Which of these lists could you add up and get something meaningful out of the
        total? \textbf{Circle} every one that works.
        \vspace{6pt}

        \hspace{1.0em}\ans{\textbf{(A)} the prices of five books in a cart}
        \hfill
        \textbf{(B)} the numbers on five players' jerseys
        \hspace{1.0em}
        \vspace{5pt}

        \hspace{1.0em}\ans{\textbf{(C)} the ages of five people in a waiting room}
        \hfill
        \textbf{(D)} the house numbers of five homes on one street
        \hspace{1.0em}
        \vspace{8pt}

        What do the ones you circled have in common that the others do not?\\[4pt]
        \ansline{A and C record an amount of something --- how many dollars, how many years.}\\[1pt]\ansline{B and D are just names written in digits, so a total means nothing.}

  \item Your school ID card has a number printed on it. Is that number \emph{measuring}
        anything about you? If not, say what it is doing instead.\\[4pt]
        \ansline{No. It measures nothing --- two students with nearby ID numbers are not}\\[1pt]\ansline{alike in any way. The number identifies me; it is a name made of digits.}

\end{enumerate}

\end{document}
